% loopkit_implementation.tex --- Section 5
% All structural details sourced from the LoopKit blog (pocoo.vaked.dev,
% 2026-06-25) and the public repo github.com/peterlodri-sec/loopkit.
\section{LoopKit Implementation}
\label{sec:loopkit}

LoopKit\footnote{\url{https://github.com/peterlodri-sec/loopkit}} is the
open-source, Python-native experimental substrate that orchestrated every
kompress training run, evaluation, and ensemble experiment in this paper. It
realizes the four-phase autonomous lifecycle and maps onto the LangChain
four-tier stacked loop taxonomy \citep{langchain2026loops}, aligning with
Osmani's five building blocks \citep{osmani2026loop} and Greyling's
operational conventions \citep{greyling2026loop}.

\subsection{The four-phase state machine}
\label{sec:loopkit:phases}

Each loop iteration executes the deterministic state machine
\textsc{Plan}~$\to$~\textsc{Execute}~$\to$~\textsc{Evaluate}~$\to$~\textsc{Decide},
persisted by an asynchronous SQLite kernel. Table~\ref{alg:fourphase}
gives the pseudocode.

\begin{table}[h]
\centering
\caption{The LoopKit four-phase iteration.}
\label{alg:fourphase}
\begin{tabular}{@{}r@{\hspace{1em}}l@{}}
\toprule
1  & \textbf{input:} hypothesis $h$, budget $B$, state $\Sigma$ \\
2  & \textbf{Phase 1: Plan} \\
3  & \quad $p \gets \mathrm{Plan}(h, \Sigma)$ \hfill {\footnotesize\textit{derive spec from $h$+$\Sigma$}} \\
4  & \textbf{Phase 2: Execute} \\
5  & \quad $r \gets \mathrm{Execute}(p)$ \hfill {\footnotesize\textit{run training/eval per spec}} \\
6  & \quad $\Sigma \gets \mathrm{SQLiteWrite}(\Sigma, p, r)$ \hfill {\footnotesize\textit{async persist}} \\
7  & \textbf{Phase 3: Evaluate} \\
8  & \quad $m \gets \mathrm{eval\_heretic}(r)$ \hfill {\footnotesize\textit{adversarial benchmark}} \\
9  & \quad $\Sigma \gets \mathrm{SQLiteWrite}(\Sigma, m)$ \\
10 & \textbf{Phase 4: Decide} \\
11 & \quad $d \gets \mathrm{Council}(m, \Sigma, B)$ \hfill {\footnotesize\textit{SHIP / RETRAIN / PIVOT}} \\
12 & \quad \textbf{if} $d = \textsc{Ship}$ \textbf{then return} $r$ \\
13 & \quad \textbf{if} $d = \textsc{Retrain}$ \textbf{then} $h \gets \mathrm{RefineHypothesis}(h, m)$ \\
14 & \quad \textbf{if} $d = \textsc{Pivot}$ \textbf{then} $h \gets \mathrm{NewHypothesis}(m, \Sigma)$ \\
15 & \textbf{goto} 2 \hfill {\footnotesize\textit{or halt if budget $B$ exhausted}} \\
\bottomrule
\end{tabular}
\end{table}

The \textsc{Decide} phase is the outer loop: the council reviews the
evaluation metric $m$ against the persistent state $\Sigma$ and remaining
budget $B$, then issues one of three verdicts. The kompress line produced 15
models through this loop---each row in the version table is one
\textsc{Plan}$\to$\textsc{Execute}$\to$\textsc{Evaluate}$\to$\textsc{Decide}
cycle, with the council (or a human acting as council) choosing the next
hypothesis.

\subsection{Mapping to the four-tier loop taxonomy}
\label{sec:loopkit:mapping}

Table~\ref{tab:loops} maps LoopKit's primitives to the LangChain four-tier
stacked loop model \citep{langchain2026loops}.

\begin{table}[h]
\centering
\caption{LoopKit instantiation of the LangChain four-tier stacked loop model.
Each tier wraps the one below it; the hill-climbing tier reaches inside and
updates the agent loop directly.}
\label{tab:loops}
\begin{tabular}{llll}
\toprule
Tier & Loop & LoopKit primitive & Role in pruning pipeline \\
\midrule
L1 & Agent        & \texttt{Loop.run()}              & training/eval iteration (Table~\ref{alg:fourphase}) \\
L2 & Verification & \texttt{evals/heretic.py}        & adversarial grader for must-keep survival \\
L3 & Event-driven & Telegram bot hook                & triggers on checkpoint events, /run /decide \\
L4 & Hill-climbing& LLM Council + trace rewriting    & decides SHIP/RETRAIN/PIVOT from eval traces \\
\bottomrule
\end{tabular}
\end{table}

\paragraph{Tier 1 (Agent loop).} \texttt{Loop.run()} implements
Table~\ref{alg:fourphase}. The abstract \texttt{Loop} class in
\texttt{loops/base.py} is subclassed per project (e.g.,
\texttt{loops/kompress/} for the pruning pipeline). Each call to
\texttt{run()} executes one full four-phase cycle.

\paragraph{Tier 2 (Verification loop).} \texttt{evals/heretic.py} is the
grader: it runs the model on adversarial prompts dense with must-keep tokens
and returns $\mathit{exact\_keep\_pct}$, $\mathit{keep\_rate}$, and
$\mathit{override\_delta}$. If the metric falls below threshold
($\mathit{exact\_keep\_pct} < 0.940$ on the 32-prompt set), the verification
loop sends the result back to Tier~1 with feedback. This is the
maker-checker separation: the training script (maker) never sees the
evaluation script's (checker) reasoning, which is why the checker catches
the silver-label ceiling that the maker could not
\citep{osmani2026loop}.

\paragraph{Tier 3 (Event-driven loop).} A Telegram bot
(\texttt{bot/main.py}) serves as the event-driven trigger. Commands
\texttt{/new}, \texttt{/run}, \texttt{/decide}, and \texttt{/history}
create, execute, decide, and query loops from chat; natural-language
messages query the council for advice. The bot reads and writes the SQLite
state on every command, so loop state survives restarts. This realizes
Osmani's automation primitive and Greyling's \texttt{loop-run-log.md}
convention in a programmatic rather than markdown form
\citep{greyling2026loop}.

\paragraph{Tier 4 (Hill-climbing loop).} The LLM Council
(\texttt{bot/council.py}) reviews evaluation traces and decides the next
action: \textsc{Ship}, \textsc{Retrain}, or \textsc{Pivot}. By default it
uses GLM-5.1, configurable to any LLM. The council's key property is that
it \emph{reaches inside} the agent loop (Tier~1) and rewrites the next
hypothesis---LangChain's insight that ``the return arrow doesn't just loop
back to the top; it reaches inside and updates the agent loop directly''
\citep{langchain2026loops}. In the kompress line, the council issued
\textsc{Retrain} three times in succession at one junction, then
\textsc{Pivot} when the metric stopped moving---the correctable-loop
invariant applied at the meta-level.

\subsection{Asynchronous SQLite state kernel}
\label{sec:loopkit:sqlite}

The state kernel (\texttt{bot/memory.py}) persists every transition of
Table~\ref{tab:loops} to an asynchronous SQLite database, serving as the
durable memory spine that Osmani identifies as the sixth building block
\citep{osmani2026loop} and that Greyling operationalizes as
\texttt{STATE.md} \citep{greyling2026loop}. Each loop iteration writes:
\begin{itemize}
  \item the hypothesis $h$ and plan $p$ (Phase~1),
  \item the raw result $r$ (Phase~2),
  \item the evaluation metrics $m$ (Phase~3),
  \item the council decision $d$ and refined hypothesis $h'$ (Phase~4).
\end{itemize}
This enables any future session to resume the loop without reconstructing
context from git history---the ``agent forgets, the repo doesn't'' principle
made concrete. The SQLite kernel is asynchronous so that long-running
training jobs on remote GPUs (vast.ai) do not block the bot's event loop;
state writes are queued and flushed when the training subprocess returns.

\subsection{Alignment with the Loop Engineering ecosystem}
\label{sec:loopkit:ecosystem}

Table~\ref{tab:ecosystem} maps LoopKit's artifacts to Greyling's operational
conventions, showing the substrate is a faithful instantiation of the
emerging standard.

\begin{table}[h]
\centering
\caption{LoopKit artifacts mapped to Greyling's Loop Engineering conventions
\citep{greyling2026loop}. Markdown conventions become programmatic
primitives; the persistence guarantee is stronger (SQLite vs.\ flat files).}
\label{tab:ecosystem}
\begin{tabular}{lll}
\toprule
Convention (Greyling) & LoopKit artifact & Enhancement \\
\midrule
\texttt{STATE.md}          & SQLite + \texttt{state.json} per loop & queryable, concurrent \\
\texttt{loop-run-log.md}   & \texttt{Experiment} dataclass + history & structured, typed \\
\texttt{loop-budget.md}    & Budget tracking in \texttt{state.json} & live, per-iteration \\
\texttt{LOOP.md}           & \texttt{GUIDE.md} & same role \\
\texttt{AGENTS.md}         & \texttt{concepts/} directory & executable reference impls \\
\texttt{loop-audit}        & Council readiness check & LLM-graded, not rule-based \\
\texttt{loop-cost}         & Budget tracking in \texttt{state.json} & live vs.\ pre-estimate \\
\bottomrule
\end{tabular}
\end{table}

\subsection{The correctable-loop invariant}
\label{sec:loopkit:correctable}

The council enforces the correctable-loop invariant: if the evaluation
metric does not improve within three iterations, halt the current direction
and pivot. Applied to the self-labeling sequence
(Section~\ref{sec:eval:C}):
\begin{itemize}
  \item v3$\to$v4: large jump ($0.942\to0.967$, $\delta$ $0.054\to0$) ---
        continue,
  \item v4$\to$v5: slight regression ($0.967\to0.961$) --- \textbf{stop}.
\end{itemize}
Two iterations, no improvement: the loop halted and pivoted to the ensemble
experiment (which revealed the Voting Ensemble Paradox). This invariant
prevents the loop from burning budget on diminishing returns---a failure
mode Osmani identifies as ``the loop changes the work, it does not delete you
from it'' \citep{osmani2026loop}.

\subsection{Automated documentation sync}
\label{sec:loopkit:docsync}

The doc-sync workflow (\texttt{.github/workflows/doc-sync.yml}) extends the
four-tier taxonomy with a continuous documentation maintenance loop. On every
push to \texttt{main}, a DeepSeek v4-flash agent reviews the diff, compares it
against the current README and author metadata, and auto-commits fixes if the
documentation has drifted from the code. The workflow always succeeds
(\texttt{continue-on-error: true}) and commits with \texttt{[skip ci]} to
avoid cascading builds---it is advisory, never blocking.

This is the maker-checker pattern applied to documentation: the code change
(maker) and the doc sync (checker) are decoupled. The checker uses a cheap
inference call (\$0.001 per push) and never blocks the shipping loop. It
realizes Greyling's \texttt{loop-run-log.md} convention
\citep{greyling2026loop} in a fully automated form: every push updates the
living documentation without human intervention, closing the loop between
code and narrative.

The doc-sync agent is part of a six-agent CI suite
(\texttt{agents/} package, invoked via \texttt{tools/ci\_agents.py} or MCP).
The suite follows a layered abstraction: each agent inherits from
\texttt{CIAgent}, returns a standardized \texttt{AgentResult}, and is
discovered via a registry. The CLI and MCP server expose the same interface,
so agents run identically in GitHub Actions, local development, or agent-to-agent
calls. Agents are advisory by design---they report pass/fail with structured
details but never block the shipping loop.

\subsection{Reproducibility and cost}
\label{sec:loopkit:repro}

All 15 kompress models were produced by LoopKit on vast.ai RTX~4090
instances at \$0.38/hr. The total compute for the v3$\to$v4$\to$v5
self-labeling chain (the paradox evidence) was \$0.55---one cold brew. The
full loopkit repository, including the kompress loop that produced these
models, is open-source at \url{https://github.com/peterlodri-sec/loopkit}
with a Colab notebook for the minimal hello-loop
(\texttt{notebooks/loopkit\_hello.ipynb}).
